\input{common/onetextpage.tex}
\newcounter{nappendix}
\renewcommand{\thenappendix}{\Alph{nappendix}}
\newcommand{\newappendix}[1]{%
    \refstepcounter{nappendix}
    \onetextpage{Appendix \thenappendix}
    \chapter*{Appendix \thenappendix: #1}
    \label{appendix:\thenappendix}
    \addcontentsline{toc}{chapter}{Appendix \thenappendix: #1}
}

%% Edit past this comment line.

\newappendix{Glossary of Terms}

\begin{tabular}{p{0.25\textwidth} p{0.7\textwidth}}
\textbf{Term} & \textbf{Definition} \\
TCCC & Thai Civil and Commercial Code (ประมวลกฎหมายแพ่งและพาณิชย์); the primary set of laws used in this simulation. \\
มาตรา (Section) & The specific numbered legal rules within the Thai code (e.g., Section 420 for Tort). \\
RAG & Retrieval-Augmented Generation; the system architecture that allows AI agents to ``look up'' Thai laws before generating a response. \\
Adversarial Multi-Agent & A setup where two AI models (Plaintiff and Defendant) are programmed with conflicting goals to simulate a courtroom battle. \\
Deka (ฎีกา) & A final ruling by the Supreme Court of Thailand, used in this project as the ``Ground Truth'' for accuracy. \\
Hallucination & When the AI generates a Thai law section or fact that does not exist in the provided database. \\
\end{tabular}

\newappendix{Agent Persona Prompts (System Instructions)}

To ensure realism, the following system prompts are used to define the behavior of the AI agents:

\section*{B.1: Plaintiff Agent Instruction}

\begin{quote}
``You are a senior Thai litigator representing the Plaintiff. Your goal is to establish liability and demand full damages. You must: (1) Use formal Thai legal language; (2) Strictly cite TCCC Sections; (3) Proactively address and dismantle the Defendant's likely arguments.''
\end{quote}

\section*{B.2: Defendant Agent Instruction}

\begin{quote}
``You are a Thai defense attorney. Your goal is to prove that your client is not liable or to minimize the damages. You must: (1) Look for procedural flaws in the Plaintiff's claim; (2) Cite conflicting Sections or exceptions in the TCCC; (3) Maintain a professional, firm tone.''
\end{quote}

\section*{B.3: Judge (Evaluator) Instruction}

\begin{quote}
``You are an impartial Legal Evaluator. Review the transcript for: (1) Legal Accuracy (Did they cite real Sections?); (2) Professionalism; and (3) Strategic Rebuttal. Provide a score from 1--5 for each category.''
\end{quote}

\newappendix{Data Schema \& Technical Logic}

\section*{C.1: Case Data Format (JSON)}

The system processes legal cases using the following structured format:

\begin{verbatim}
{
  "case_id": "CIVIL-2026-001",
  "domain": "Breach of Contract",
  "applicable_law": ["Section 213", "Section 387"],
  "facts": {
    "plaintiff_evidence": "Evidence of payment and non-delivery of goods.",
    "defendant_evidence": "Force Majeure due to extreme weather conditions."
  },
  "ground_truth_verdict": "Plaintiff Wins (Partial Damages)"
}
\end{verbatim}

\section*{C.2: Simulation Workflow}

The system follows a four-step execution loop:

\begin{enumerate}
    \item Selection: User picks a Deka case or enters facts.
    \item Retrieval: The RAG system queries the Vector DB for matching TCCC Sections.
    \item Debate: Agents exchange turns (Opening $\rightarrow$ Rebuttal $\rightarrow$ Closing).
    \item Evaluation: The Judge AI compares the outcome against the actual Supreme Court verdict.
\end{enumerate}

\newappendix{Software Tools \& Libraries}

\begin{itemize}
    \item Core LLM: ...
    \item Thai NLP: PyThaiNLP (for segmenting Thai legal text).
    \item Vector Database: ... (storing the Thai Civil and Commercial Code).
    \item Framework: .... (managing the RAG pipeline).
    \item Datasets: Thai-Law-Dataset (PyThaiNLP) and NitiBench.
\end{itemize}

(Everything in here is a placeholder. We will fill in the actual tools and libraries we use in the final version of the SRS document.)
